\section{Background and related work}
\label{sec:background}

\subsection{Typed decision models}
\label{sec:tdm}

A typed decision model receives a \emph{state} (text or structured data) and one or more questions, each with an option set declared by the caller, and for each question returns a probability distribution over exactly those options without generating free text~\cite{arx23886,arx26758,arx32160}.
\jev{} offers three question types: \emph{Choice} selects among listed options, \emph{Score} rates the state against ordered rubric levels, and \emph{Noul} returns the probability that a statement is true~\cite{typesafe_intro,arx32160}.
The hosted model is closed. TypeSafe describes training by \emph{reinforcement learning for calibrated decisions} (RLCD) but publishes neither the architecture, the reward, the data nor the weights~\cite{almeida2026introducing,typesafe_primer}.
Open-weight models with the same interface, which read the hidden state against the option labels in a decoder or read option markers or definitions in an encoder~\cite{arx23886,arx23959,arx26758}, make the mechanism inspectable.
Masking options outside the declared set makes every answer a valid option, which guarantees form but not correctness~\cite{arx26758}.
Hosted \jev{} is not deterministic~\cite{arx32160,arx01006}.
The vendor lists \$0.042 per million input tokens with free output~\cite{typesafe_models}, and its launch post cites 70--500\,ms per call~\cite{almeida2026introducing}.
Reported cost reductions against generative evaluators range from a few-fold to two orders of magnitude, depending on the comparator and the pricing assumed~\cite{arx29429,arx27535}, and the advantage is not universal~\cite{arx23136,arx00437}.

\subsection{Related work}
\label{sec:related}

The questions raised by typed safeguards are new instances of older problems.
Platform moderation has long combined automated classifiers with human review, and its critics have shown that automation hides contestable policy choices inside technical systems~\cite{gorwa2020algorithmic,gillespie2018custodians}.
Safety classifiers that read token probabilities~\cite{inan2023llama,markov2023holistic,sharma2025constitutional} and programmable guardrails~\cite{rebedea2023nemo} are the closest technical precursors of a typed safety valve.
Prompt injection~\cite{perez2022ignore,greshake2023not} and adversarial attacks~\cite{zou2023universal,toyer2024tensor} are the corresponding threats.
Automated hate-speech detection shows racial and dialect bias~\cite{sap2019risk,davidson2019racial} and fails targeted functional tests~\cite{rottger2021hatecheck}.
Aggregating annotators' labels by majority vote discards the disagreement a safeguard needs~\cite{davani2022dealing,plank2022problem}.
Chinese offensive-language benchmarks exist~\cite{deng2022cold,lu2023toxicn}, but none labels the constructs of \pol{}.
Inferring emotion from observable signals rests on weak scientific foundations~\cite{barrett2019emotional,stark2021ethics}.
Calibration degrades under distribution shift~\cite{guo2017calibration,ovadia2019trust} and can hold on average while failing for subgroups~\cite{hebertjohnson2018multicalibration}.
Abstention, learning to defer and conformal risk control formalise when a model should hand a case on~\cite{chow1970optimum,geifman2017selective,madras2018predict,angelopoulos2024conformal}.
Language models are sensitive to option order, label names and prompt format~\cite{zhao2021calibrate,pezeshkpour2024large,sclar2024quantifying}, and LLM evaluators show position bias and correlated errors~\cite{wang2024fair,goel2025great}, for which panels and cascades are the established remedies~\cite{verga2024replacing,chen2023frugalgpt}.
Human oversight of automated decisions often fails through over-reliance and misplaced responsibility~\cite{parasuraman2010complacency,green2022flaws,elish2019moral}, and automated adjudication threatens notice, reasons and the opportunity to contest~\cite{citron2008technological}, which data-protection, platform and AI law now partly codify~\cite{gdpr2016,dsa2022,euaiact2024,santaclara2021}.

\subsection{\pol{} among AI-governance paradigms}
\label{sec:paradigms}

\oa{D} compares \pol{}, as a text, with risk-based regulation, platform content governance, training-time alignment to a written constitution or spec, inference-time guardrails and decentralised governance, five paradigms under which automated judgement is already proposed or deployed.
Most of them specify actors and procedures rather than the judgement itself.
Risk-based regulation assigns duties by risk tier.
The EU instruments prescribe no interception, while China's Interim Measures require providers to stop generating or transmitting unlawful content once found (Art.~14(1)), without specifying a mechanism~\cite{euaiact2024,cac2023genai}.
Platform governance specifies notice, reasons and appeal around decisions whose criteria remain the platform's own policies~\cite{dsa2022,santaclara2021}.
Training-time constitutions and specs state norms in natural language but implement them inside one vendor's model, and no output records which principle it followed~\cite{bai2022constitutional,openai2025modelspec}.
Guardrails specify the intercept precisely, as code, but state no duty about reasons, records or the human role~\cite{inan2023llama,rebedea2023nemo}.
DAOs specify how rules execute and who votes, but not how content or behaviour is judged~\cite{buterin2014ethereum,hassan2021decentralized}.

Five properties characterise \pol{} in this comparison, each with a cost.
(i) It places a model-agnostic intercept inside existing LLMs (\clause{5.4}, \clause{5.5}) but leaves its mechanism, threshold and appeal unspecified.
(ii) Its three-valued construct, Love Language, hate language and the state of absence, is declared not to be a set of moral labels (\clause{1.5}, \clause{4.3.2}) and has not been operationalised beyond 20 pilot items~\cite{polgov}.
(iii) The text is open, versioned and citable, but it is controlled by its originator and has not been peer reviewed.
(iv) Auditability duties for welfare decisions, a verifiable decision chain, explanation on request and access to all data and logs (\art{5}(5), \art{9}(2)--(3)), conflict with the privacy of the persons judged and with keeping decision probabilities from attackers.
(v) Population-level aims appear in citable clauses: the text permits per-person quantification and provides for AI-only welfare decisions (\clause{5.6}, which says governed AI ``can'', not must, do the former; \art{2}(1), \art{9}(7)).
AI-only welfare decisions sit against EU human-oversight provisions (directly for the systems they cover, by analogy otherwise), and per-person quantification resembles the social scoring that the AI Act prohibits.
\pol{} and the Interim Measures are the two texts compared that join, in one instrument, an inference-time intercept, the construct it applies and duties attached to decisions.
Both leave a gap at the individual block, since neither requires that a person whose content is blocked be told of the block.
We study \pol{} for these properties, not for uniqueness.
